\section{Method}
\label{sec:method}

% DRAFT (Part T, 2026-11-27). Written from the code on branch overnight-nov26:
% rua_harness/command.py, policy/{model,variants,data,label_commands,
% commands_gt,run_skill,train}.py, policy/commander/. Every design choice that
% has no result yet is marked with \draftnote{no result yet}.

\draftnote{DRAFT. Describes what is built. No trained-policy result exists yet.}

An LLM agent decides what to do. A small learned policy decides how to move.
Between them sits one \emph{command}: the message the agent sends each time it
hands the policy a job. We fix the command's content and vary only its
\emph{format}, the way that content is shown to the policy. For each format we
retrain the same policy on the same demonstrations.

\subsection{The command}
\label{sec:command}

A command names a skill and the things it acts on. The skills are
\texttt{pick\_place}, \texttt{open}, \texttt{close}, \texttt{slide},
\texttt{press}, \texttt{turn}, \texttt{insert} and \texttt{navigate}. Each
skill has required fields (for example, \texttt{pick\_place} needs an
\texttt{object} and a \texttt{target}) and optional ones (a \texttt{source}
name, or a \texttt{direction} such as \texttt{on} or \texttt{out}). The
\texttt{object} and \texttt{target} are references: a short name and,
optionally, a 3D point in the robot base frame where the gripper should act.

The command carries two kinds of content, and we keep them apart on purpose.
The \emph{semantic channel} is the names: skill, object, source, target,
direction. It is text and goes to the policy's frozen language encoder. The
\emph{geometric channel} is the points (or, in one format, a pose path). It
goes to a small learned encoder, or is drawn on the images. Numbers never pass
through the language encoder, and names never pass through the point encoder.

\subsection{Five command formats}
\label{sec:formats}

All formats share the same demonstrations, cameras (left and right agent
views), state, action chunk and training budget. They differ only in how the
command enters the policy.

\paragraph{\texttt{sentence}.} The original RoboCasa instruction, for example
``Open the left drawer.'' This is stock \dem{} and the control for every
other format.

\paragraph{\texttt{kv\_text}.} Structured text with a fixed field order, no
punctuation and no numbers, for example
\texttt{skill pick place object apple from counter to cabinet}. It carries the
same names as the sentence in a regular form that an agent can fill in
without writing prose. The formats below all keep this text and add geometry
to it, so each is compared against \texttt{kv\_text} as well as against
\texttt{sentence}.

\paragraph{\texttt{points\_2d}.} The grasp point and the place point are
projected into each agent-view camera and drawn on its image as filled
circles (radius 4 pixels at $256\times256$, green for grasp, red for place,
each with a dark outline). Points behind the camera or outside the image are
not drawn. The policy architecture is unchanged; it reads the points from the
pixels.

\paragraph{\texttt{points\_3d} and \texttt{points\_3d\_rgb}.} The two 3D
points, in the robot base frame of the current frame, become two extra tokens.
Each point goes through random Fourier features (16 frequencies, scale 10,
concatenated with the raw coordinates) and a three-layer MLP (hidden width
256, GELU) to one token of the language width, plus a learned embedding for
its slot (grasp or place). The two tokens are appended to the 32 language
tokens, with a mask entry set to false when a point is absent. The \dem{}
head is unchanged and sees $32+2$ tokens in its language stream. The plan
gives \texttt{points\_3d} a point-map image input (per-pixel 3D coordinates
from rendered depth) in place of color, and \texttt{points\_3d\_rgb} keeps
color, so the effect of the point tokens and of the image input can be
separated. \draftnote{In the current code both use color images. The point-map
input is not yet wired into the training loader; only
\texttt{points\_3d\_rgb} is queued. No result yet.}

\paragraph{\texttt{pose\_path}.} Eight future gripper poses, one every 8
frames (0.4\,s at 20\,Hz), each a position and a quaternion expressed in the
current base frame. The same kind of encoder (input size 7, eight slots)
turns them into eight tokens appended to the language stream. Near the end of
a demonstration the last pose repeats. This format tells the policy how to
move, not only where, and it is the hardest for an agent to write. At test
time a perfect path is available only from a demonstration's own future, so
this format is scored on demonstration scenes only. \draftnote{Planned with a
point-map input; the current code uses color. No result yet.}

\subsection{Hindsight command labels}
\label{sec:labels}

The demonstrations come with an instruction but no points. We label each
demonstration with the command it should have received, in hindsight, from
its own recorded actions. The gripper action is $-1$ (open) or $+1$
(closed). The first frame where the gripper closes gives the grasp, and the
next frame where it opens gives the place. The gripper position at those
frames, computed from the state vector (base pose plus end effector relative
to the base), becomes the object point and the target point. Names come from
the task and the episode metadata (object category, fixture names).

Some demonstrations never close the gripper (buttons, knobs, levers, and
pushed doors). For these we use the gripper position at the frame the
demonstration first succeeds. Demonstrations end 17 frames after first
success, so that frame is the length minus 17. \texttt{NavigateKitchen}
uses the final base position as the target. Each label records which method
produced it; the counts per task are in \cref{sec:setup}.

Labels are stored in the world frame. At training frame $t$ they are moved
into the robot base frame of frame $t$, using the base pose in the state
vector, so the command is always relative to where the robot is now.
\draftnote{The success-frame labels are not yet checked against the true
fixture part.}

\subsection{Noise on the command during training}
\label{sec:noise}

An agent's points will be off by a few centimeters. We add Gaussian noise
during training: each coordinate of each point gets noise with a standard
deviation of 2\,cm, and in \texttt{pose\_path} each pose also gets a random
rotation with a standard deviation of $3^\circ$ per axis. A fresh random
generator is drawn per sample, so data-loader workers do not share noise.
\draftnote{The noise scale is a design choice with no result yet. A sweep (0,
1, 3, 5\,cm) is in the backlog.}

\subsection{The policy}
\label{sec:policy}

The policy is \dem{} \citep{sun2026dem}. It encodes two camera images at
$256\times256$ with a fine-tuned DINOv3 ConvNeXt-B \citep{simeoni2025dinov3},
encodes up to 32 language tokens with a frozen NeoBERT
\citep{lebreton2025neobert}, and reads a 16-number robot state (base position
and orientation, end-effector position and orientation relative to the base,
two gripper joints). A MeanFlow action head \citep{geng2025meanflow}
predicts a chunk of 16 actions of 12 numbers each (base motion and torso,
control mode, end-effector position and rotation deltas, gripper) in one
step, and the robot executes the first 8 before predicting again. We do not
change the backbone or the head. The only additions are the geometric
encoder of \cref{sec:formats} and, for \texttt{points\_2d}, the drawn
markers.

\subsection{The agent side}
\label{sec:agent}

The agent is a coding-agent LLM that we call the \emph{commander}. It
controls one episode through two tools. \texttt{look} returns the current
left, right and wrist camera images. \texttt{run\_skill} takes a command in
the form of \cref{sec:command}, except that each reference gives a camera
and a pixel $(u,v)$ instead of a 3D point. The harness lifts the pixel to 3D
using that camera's rendered depth, intrinsics and pose, moves the point into
the robot base frame, and runs the policy until the skill stops or its time
runs out. It then returns a status and new images, never simulator ground
truth. A third call, \texttt{done}, ends the episode. Robot time across all
skills is capped at the task's horizon.

The policy has no learned ``done'' output, so a skill stops when the
simulator's success check turns true, when its time limit runs out, or at a
step limit. The success flag is used only to stop the skill and is not shown
to the agent. \draftnote{This leaks a little ground truth (a skill ends early
when it succeeds). A learned done output is planned.}

Simulator ground truth is used for grading only. The one exception is the
``perfect command'' used to score formats without an agent
(\cref{sec:setup}), which plays the role of an ideal agent.
